\documentclass[11pt, letterpaper]{article}
\usepackage[utf8]{inputenc}
\usepackage[T1]{fontenc}
\usepackage[spanish]{babel}
\usepackage[margin=2.5cm]{geometry}
\usepackage{amsmath}
\usepackage{booktabs}
\usepackage{array}
\usepackage{enumitem}
\usepackage{parskip}
\usepackage{titlesec}
\usepackage{fancyhdr}
\usepackage{listings}
\usepackage{xcolor}
\usepackage{hyperref}
\usepackage{caption}
\usepackage{float}
\usepackage{tikz}
\usetikzlibrary{shapes.geometric, arrows.meta, positioning, fit, backgrounds}

% ── Estilo de código ──────────────────────────────────────────────────────────
\lstset{
  basicstyle=\ttfamily\small,
  backgroundcolor=\color{gray!8},
  frame=single,
  rulecolor=\color{gray!40},
  breaklines=true,
  numbers=none,
  keywordstyle=\color{blue!70!black},
  commentstyle=\color{green!50!black},
  stringstyle=\color{orange!80!black},
}

% ── Encabezado ────────────────────────────────────────────────────────────────
\pagestyle{fancy}
\fancyhf{}
\rhead{IIC 2440 --- Tarea 1}
\lhead{Parte 1: Data Warehouse}
\rfoot{\thepage}

\titleformat{\section}{\large\bfseries}{}{0em}{}[\titlerule]
\titleformat{\subsection}{\normalsize\bfseries}{}{0em}{}

% ─────────────────────────────────────────────────────────────────────────────
\begin{document}

\begin{center}
  {\Large \textbf{Tarea 1 --- IIC 2440: Procesamiento de Datos Masivos}} \\[4pt]
  {\large Parte 1: Data Warehouse y Pipeline ETL} \\[2pt]
  {\small Pontificia Universidad Católica de Chile \quad --- \quad Primer Semestre 2026}
\end{center}

\vspace{0.5em}
\hrule
\vspace{1em}

% ═══════════════════════════════════════════════════════════════════════════════
\section{Esquema Estrella}

El warehouse sigue un esquema estrella con una tabla de hechos central
(\texttt{fact\_news}) y tres tablas de dimensiones. Cada artículo en
\texttt{fact\_news} reemplaza sus atributos categóricos por llaves foráneas
hacia las dimensiones correspondientes, eliminando redundancia y facilitando
consultas analíticas.

\subsection{Tabla de hechos: \texttt{fact\_news}}

\begin{table}[H]
\centering
\begin{tabular}{lll}
\toprule
\textbf{Columna} & \textbf{Tipo} & \textbf{Descripción} \\
\midrule
\texttt{article\_id}       & string  & Identificador único del artículo (hash) \\
\texttt{date\_id}          & int32   & FK $\to$ \texttt{dim\_date} \\
\texttt{source\_id}        & int32   & FK $\to$ \texttt{dim\_source} \\
\texttt{region\_id}        & int32   & FK $\to$ \texttt{dim\_region} \\
\texttt{title}             & string  & Título del artículo (encoding corregido) \\
\texttt{body}              & string  & Cuerpo del artículo (encoding corregido) \\
\texttt{title\_word\_count} & int32  & Cantidad de palabras en el título \\
\texttt{body\_word\_count}  & int32  & Cantidad de palabras en el cuerpo \\
\bottomrule
\end{tabular}
\caption{Esquema de \texttt{fact\_news}.}
\end{table}

La tabla está \textbf{particionada por año y mes} (\texttt{year=YYYY/month=MM}),
siguiendo la convención Hive. Cada partición se almacena en un archivo Parquet
independiente (\texttt{part-0.parquet}), lo que permite leer solo el rango de
fechas necesario en consultas posteriores.

\subsection{Tablas de dimensiones}

\begin{table}[H]
\centering
\begin{tabular}{lll}
\toprule
\textbf{Tabla} & \textbf{Clave primaria} & \textbf{Atributos} \\
\midrule
\texttt{dim\_date}   & \texttt{date\_id}   & fecha, año, mes, día, día de semana, semana del año \\
\texttt{dim\_source} & \texttt{source\_id} & nombre del dominio (normalizado a minúsculas) \\
\texttt{dim\_region} & \texttt{region\_id} & nombre, número romano oficial, abreviación \\
\bottomrule
\end{tabular}
\caption{Tablas de dimensiones.}
\end{table}

\texttt{dim\_region} incluye una fila especial \texttt{Desconocida} (id = 17)
para artículos sin mención detectable de región, garantizando que
\textbf{todos} los registros de \texttt{fact\_news} tengan una región válida
y la integridad referencial no se rompa.

% ═══════════════════════════════════════════════════════════════════════════════
\section{Pipeline ETL}

El pipeline se implementó en Python puro (pandas + pyarrow) y se puede
ejecutar con un solo comando: \texttt{python main.py}. Consta de cuatro
pasos secuenciales.

\begin{figure}[H]
\centering
\begin{tikzpicture}[
  node distance=0.6cm and 1.2cm,
  box/.style={rectangle, draw, rounded corners=3pt,
              minimum width=2.8cm, minimum height=0.85cm,
              align=center, font=\small},
  arr/.style={-{Stealth}, thick},
]
  \node[box, fill=blue!10]   (csv)    {CSV crudo};
  \node[box, fill=orange!15, right=of csv]  (scan)   {Paso 1\\Escaneo};
  \node[box, fill=green!15,  right=of scan] (dims)   {Paso 2\\Dimensiones};
  \node[box, fill=yellow!20, right=of dims] (fact)   {Paso 3\\fact\_news};
  \node[box, fill=red!10,    right=of fact] (val)    {Paso 4\\Validaciones};

  \draw[arr] (csv)  -- (scan);
  \draw[arr] (scan) -- (dims);
  \draw[arr] (dims) -- (fact);
  \draw[arr] (fact) -- (val);
\end{tikzpicture}
\caption{Flujo del pipeline ETL.}
\end{figure}

\subsection{Paso 1 --- Escaneo ligero (\textit{pass 1})}

Se lee el CSV usando únicamente las columnas \texttt{publish\_date} y
\texttt{source} (sin cargar \texttt{body}), en chunks de 100\,000 filas.
El objetivo es construir dos diccionarios:
\begin{itemize}[noitemsep]
  \item \texttt{unique\_dates}: \texttt{fecha\_str} $\to$ \texttt{date\_id}
  \item \texttt{unique\_sources}: \texttt{dominio} $\to$ \texttt{source\_id}
\end{itemize}
Esta pasada es rápida porque evita deserializar los campos de texto largos.
Al finalizar se conoce también el conteo total de filas crudas
(\texttt{raw\_row\_count}), usado en las validaciones.

\subsection{Paso 2 --- Construcción de dimensiones}

Con los diccionarios del paso anterior se construyen las tres dimensiones
en memoria y se persisten como Parquet:

\begin{itemize}[noitemsep]
  \item \textbf{\texttt{dim\_date}}: se parsea cada fecha única con
    \texttt{datetime.strptime} para extraer año, mes, día, día de la semana
    e ISO-semana.
  \item \textbf{\texttt{dim\_source}}: se normaliza el dominio a minúsculas
    para evitar duplicados por capitalización.
  \item \textbf{\texttt{dim\_region}}: se construye estáticamente desde la
    tabla de aliases definida en \texttt{regions.py}.
\end{itemize}

\subsection{Paso 3 --- ETL principal (\textit{pass 2})}

Se recorre el CSV nuevamente en chunks. Por cada chunk:

\begin{enumerate}[noitemsep]
  \item \textbf{Limpieza}: se eliminan filas sin \texttt{article\_id} o
    \texttt{publish\_date}, y duplicados por \texttt{article\_id}.
  \item \textbf{Corrección de encoding}: se aplica \texttt{fix\_encoding}
    sobre \texttt{title} y \texttt{body} (ver Sección~\ref{sec:encoding}).
  \item \textbf{Asignación de llaves foráneas}: se hace \textit{map} de
    \texttt{publish\_date} $\to$ \texttt{date\_id} y de \texttt{source}
    $\to$ \texttt{source\_id} usando los diccionarios del paso~1.
  \item \textbf{Inferencia de región}: se aplica \texttt{assign\_regions}
    (ver Sección~\ref{sec:region}) y se asigna \texttt{region\_id}.
  \item \textbf{Word count}: se cuenta el número de palabras en
    \texttt{title} y \texttt{body} de forma vectorizada.
  \item \textbf{Escritura particionada}: se agrupa el chunk por
    \texttt{(year, month)} y cada grupo se escribe en su \texttt{ParquetWriter}
    correspondiente.
\end{enumerate}

Los \texttt{ParquetWriter}s se mantienen abiertos en un diccionario
\texttt{writers} durante todo el procesamiento. Dado que el CSV no está
ordenado cronológicamente, artículos de la misma partición pueden aparecer en
chunks distintos; mantener el writer abierto permite acumular esas filas en
el mismo archivo sin reescribirlo. Al terminar (o si ocurre un error),
todos los writers se cierran en un bloque \texttt{finally}, lo que
materializa el \textit{footer} de Parquet y deja los archivos válidos.

\subsection{Paso 4 --- Validaciones}

Se describen en la Sección~\ref{sec:validaciones}.

% ═══════════════════════════════════════════════════════════════════════════════
\section{Decisiones de diseño}

\subsection{Procesamiento por chunks}
\label{sec:chunks}

El dataset completo ocupa $\approx$4\,GB en disco. Cargarlo en memoria de
una sola vez requeriría varios GB de RAM adicionales por las estructuras
intermedias de pandas. En cambio, procesar en chunks de 100\,000 filas
mantiene el uso de memoria acotado independientemente del tamaño del archivo,
a costa de dos pasadas en lugar de una.

\subsection{Corrección de encoding}
\label{sec:encoding}

El dataset presenta dos tipos de corrupción de texto:

\begin{itemize}[noitemsep]
  \item \textbf{Mojibake clásico} (ej. \texttt{Ríos} en lugar de
    \textit{Ríos}): causado por leer bytes UTF-8 como Latin-1. Se corrige
    con la librería \texttt{ftfy}.
  \item \textbf{Sustituciones CJK} (ej. \texttt{U+8C69} en lugar de
    \textit{á}): artefacto específico de este dataset donde vocales
    acentuadas fueron reemplazadas por caracteres CJK. Se corrige con un
    diccionario de mapeo manual.
\end{itemize}

Para evitar el costo de llamar a \texttt{ftfy} en cada fila, se aplica
primero una búsqueda con expresión regular que detecta bytes sospechosos
(\texttt{\textbackslash xc3}, \texttt{\textbackslash xc2}, caracteres CJK).
Solo si hay una coincidencia se invoca \texttt{ftfy}.

\subsection{Inferencia de región}
\label{sec:region}

Se definió una tabla de aliases por región (topónimos, ciudades principales,
número romano oficial) con más de 100 términos en total. La inferencia opera
en tres etapas vectorizadas sobre cada chunk:

\begin{enumerate}[noitemsep]
  \item \textbf{Normalización}: se convierte el texto a minúsculas y se
    eliminan tildes (NFD + ASCII), para que las búsquedas sean
    insensibles a mayúsculas y acentos.
  \item \textbf{Puntuación}: se cuentan las menciones de cada alias con
    \texttt{str.count} y se aplica la ponderación:
    \[
      \text{score}(r) = 3 \cdot \text{hits\_title}(r)
                      + 1 \cdot \text{hits\_body}(r)
    \]
    donde \texttt{body} se trunca a 800 caracteres para reducir el cómputo.
  \item \textbf{Desambiguación}: la región con mayor puntaje gana.
    Si el puntaje máximo es 0 o hay empate, se asigna \texttt{Desconocida}.
\end{enumerate}

La ponderación $3\!:\!1$ refleja que las menciones en el título son más
informativas: el título suele nombrar explícitamente el lugar de la noticia,
mientras que el cuerpo puede mencionar otras regiones de forma contextual.

\subsection{Formato de almacenamiento}

Se eligió \textbf{Parquet} (columnar, comprimido) sobre CSV o JSON por tres
razones: (1) ocupa significativamente menos espacio en disco; (2) permite
leer solo las columnas necesarias en cada consulta; (3) preserva los tipos
de datos sin necesidad de parseo adicional. La partición Hive
(\texttt{year=/month=}) permite además que herramientas como PyArrow o
DuckDB omitan particiones completas al filtrar por fecha.

% ═══════════════════════════════════════════════════════════════════════════════
\section{Validaciones}
\label{sec:validaciones}

Se implementaron cinco validaciones que se ejecutan automáticamente al
finalizar el pipeline. Para eficiencia, \texttt{fact\_news} se carga con
solo las columnas de llaves (\texttt{article\_id}, \texttt{date\_id},
\texttt{source\_id}, \texttt{region\_id}).

\begin{table}[H]
\centering
\begin{tabular}{clp{8cm}}
\toprule
\textbf{\#} & \textbf{Nombre} & \textbf{Qué verifica} \\
\midrule
1 & Consistencia referencial &
  Ninguna FK en \texttt{fact\_news} apunta a un ID inexistente en
  \texttt{dim\_date}, \texttt{dim\_source} o \texttt{dim\_region}. \\
2 & Conteo de filas &
  El número de filas en \texttt{fact\_news} no supera el número de filas
  en el CSV crudo (las diferencias se deben a filas descartadas). \\
3 & PKs sin duplicados &
  Cada tabla de dimensiones tiene valores únicos en su clave primaria. \\
4 & Particiones correctas &
  Para cada par (año, mes) presente en \texttt{fact\_news}, existe el
  archivo \texttt{part-0.parquet} en la carpeta de partición correspondiente. \\
5 & \texttt{article\_id} únicos &
  No existen \texttt{article\_id} duplicados en \texttt{fact\_news},
  garantizando la unicidad de la clave natural. \\
\bottomrule
\end{tabular}
\caption{Resumen de validaciones implementadas.}
\end{table}

% ═══════════════════════════════════════════════════════════════════════════════
\section{Estructura del warehouse}

\begin{lstlisting}
warehouse_data/
    dim_date/
        dim_date.parquet
    dim_source/
        dim_source.parquet
    dim_region/
        dim_region.parquet
    fact_news/
        year=2023/
            month=01/part-0.parquet
            month=02/part-0.parquet
            ...
        year=2024/ ...
        year=2025/ ...
\end{lstlisting}

El pipeline completo se ejecuta con:
\begin{lstlisting}[language=bash]
python main.py                   # dataset completo
python main.py --sample          # noti.csv (50k filas, para pruebas)
python main.py --chunk-size 200000  # chunk mas grande
\end{lstlisting}

\end{document}
